% Pursuing Stacks (cotes 134-2 to 134-8): archive cote 134-6, site page 9 (PDF page 10), typescript page 328.
% typescript. Draft compared: ps5.tex:1778-1852.
% AI-transcribed by Claude Opus 5.5 (claude-opus-5-5, adaptive thinking, effort high), Message Batches API, batch ps3 (msgbatch_01XdFmuVVBAcgc7EFxdtCqsj), 2026-11-12; not yet checked by eye.
\scanpage{134-6/9}
\mspage{328}

(18) \qquad $A^{\wedge}_{ab} \longrightarrow \hand{\Delta}^{\wedge}_{ab}$

we obtain by composing the two arrows in the right hand vertical row,
transforms weak equivalences into same, which is needed in order to
deduce \xout{by} (14) by passing to the localized categories. When the two
assumptions are satisfied, then the functor (14) obtained from (18)
does factorize (6) as required, and it should be clear too that it
satisfies the compatibility (15).

Thus it seems there are good prospects for getting canonical
functors
\xout{(19)} \qquad $\mathrm{Hot}_{A} \longrightarrow \mathrm{Hotab}_{A} \longrightarrow \mathrm{Hotab}$ \quad ,
whose composition is (6), i.e. inserting into the commutative diagram
\diagram{\[
(19)\qquad
\begin{tikzcd}
\mathrm{Hot}_{A} \arrow[r] \arrow[d] & \mathrm{Hotab}_{A} \arrow[d] \\
\mathrm{Hot} \arrow[r] & \mathrm{Hotab}
\end{tikzcd}
\quad .
\]}
The next question then which arises is \uline{whether the second vertical}
\uline{arrow} (namely (14)) \uline{is an equivalence}, \uline{whenever the first one is}, i.e.
when $A$ is a pseudo test category, or whether this is true if we make
some familiar extra assumption on $A$, \inserted{such} as being an actual test category\inserted{say}.

I feel I am getting gradually back into thin air conjecturing,
I would'nt go on too long this way\hand{!} This \xout{Hxtab} whole $\mathrm{Hotab}_{A}$ business
was just a digression, which then took me longer than expected, it
does'nt seem to have much to do with what I have been out for\unsure{*} in
this section, namely afterthoughts about \uline{boundary operations} \uline{in a test}
\uline{category}, which are designed to give a ``computational'' description
of the canonical functor (6), \inserted[typed]{the latter being} defined without any restriction nor dif-%
ficulty for \uline{any} small category $A$. More accurately still , we want to
describe the composition of (6) with the canonical functor $A^{\wedge} \to \mathrm{Hot}_{A}$ ,
namely $A^{\wedge} \longrightarrow \mathrm{Hotab}$ , via a suitable functor $A^{\wedge} \longrightarrow \mathrm{Ch.(Ab)}$, with
the expectation that the latter should factor through $A^{\wedge}_{ab}$ via the
abelianization functor $\mathrm{Wh}_{A}$ . In other words, we \xout{wantxto} are looking
for commutative diagrams
\diagram{\[
(20)\qquad
\begin{tikzcd}
 & A^{\wedge} \arrow[dl] \arrow[rr] \arrow[dd, "\hand{K}"] & & \mathrm{Hot} \arrow[dd] \\
A^{\wedge}_{ab} \arrow[dr, "\hand{L}"'] & & & \\
 & \mathrm{Ch.(Ab)} \arrow[rr] & & \mathrm{Hotab}
\end{tikzcd}
\quad ,
\]}

\margin{9}
